\scanpage{64}\mspage{63}
ingrédients principaux pour décrire l'atelier, via $(M,\leq,\ll,\mathrel{\tikz[baseline=-.45ex,x=1em,y=1em]{\draw[line width=.055em] (.52,.27)--(.05,.055)--(.48,-.11); \draw[line width=.055em] (.14,.21)--(.46,.04)--(0,-.17);}})$) en termes
ensemblistes, au moyen de la seule
relation $\mathrel{\tikz[baseline=-.45ex,x=1em,y=1em]{\draw[line width=.055em] (.52,.27)--(.05,.055)--(.48,-.11); \draw[line width=.055em] (.14,.21)--(.46,.04)--(0,-.17);}}$ dans $L$ (puisque
\deleted{\unsure{si|il}} il est clair que pour deux \uline{lieux},
\begingroup\makeatletter\tagsleft@true\makeatother
\[
x\parallel y\quad\Longleftrightarrow\quad x \mathrel{\tikz[baseline=-.45ex,x=1em,y=1em]{\draw[line width=.055em] (.52,.27)--(.05,.055)--(.48,-.11); \draw[line width=.055em] (.14,.21)--(.46,.04)--(0,-.17);}} y\quad\text{et}\quad x\ne y
\tag*{(1.59)}
\]
\endgroup

\deleted{Le seul}
Cela mettrait donc un poids particulier
sur cette relation de $\underset{\unsure{L|\alpha}}{\mathrel{\tikz[baseline=-.45ex,x=1em,y=1em]{\draw[line width=.055em] (.52,.27)--(.05,.055)--(.48,-.11); \draw[line width=.055em] (.14,.21)--(.46,.04)--(0,-.17);}}}$ (ou ce qui
revient au $=\quad\underset{L}{\parallel}$) \uline{\deleted{\unsure{de|dans}}\unsure{dans}} $L$. Cet
\unsure{axiome} est bien raisonnable, y
compris dans des cas non ``ensembliste'',
mais le seul hic, c'est que il n'est
\uline{pas} satisfait dans les cas simpliciaux
les plus importants, cf exemples page \ldots.
Mais dans ce cas, on peut se rabattre sur
une \deleted{\unsure{propriét|propriété}} propriété qui fera le $=$
usage

% slanted
\margin{Variante ``simpliciale'' de \uline{At} 9$'$}
\uline{At \deleted{\unsure{fac|fer}}} \uline{9}$''$ Soient $X,Y\in M$ et $F$ comme
dans \unsure{Atfer|Atfac} 1. Alors
\[
X \mathrel{\tikz[baseline=-.45ex,x=1em,y=1em]{\draw[line width=.055em] (.52,.27)--(.05,.055)--(.48,-.11); \draw[line width=.055em] (.14,.21)--(.46,.04)--(0,-.17);}} Y\ \Longleftrightarrow\
\left\{\begin{aligned}
&\forall x\in\operatorname{omb}X\setminus\operatorname{omb}F\quad\text{et}\quad y\in\operatorname{om}Y\setminus\operatorname{omb}F,\\
&\hspace{9em}\text{on a }\ x\parallel y\\
&F\ \text{est \inserted{vide ou est} ``élémentaire'' (i.e. }F\in M\\
&\qquad\text{\unsure{(i.e.|i.e.} c'est une multistrate commune).}
\end{aligned}\right.
\]
\margin{\tikz[baseline=(note.base)]{\node[inner sep=2pt,outer sep=0pt,text width=\dimexpr\marginparwidth-4pt\relax,align=left] (note) {\uline{NB} On peut donner \unsure{variantes|variante} avec toute autre condition, sur l'ens ordonné $\widetilde F$.}; \draw[rounded corners=4pt] (note.north west)--(note.north east)--(note.south east);}}
